\section{The F-PSU Protocol}
\label{sec:protocol}

	
\subsection{System Model and Threat Model}
% ============================================================

\textbf{System model.} Two parties $P_0, P_1$ hold private sets $X_0, Y_0
\subseteq \{0,1\}^\sigma$. The protocol runs over $T = \poly(\kappa)$ epochs.
At epoch $t$, each party provides a delta: elements to add ($\Delta^+$) and
delete ($\Delta^-$), disjoint from each other and from the current set. After
epoch $t$, both parties learn the updated union $X_{\cup,t}$.

\textbf{Updatable efficiency.} A protocol is \emph{$1$-updatable} if per-epoch
communication $|\mathsf{trans}_t| = O(|\Delta_t| \cdot \poly(\sigma, \kappa))$,
scaling with delta size rather than total set size.

\textbf{Threat model.} Semi-honest adversary with adaptive corruption of one
party at epoch $\tau$ (corrupting both simultaneously voids security guarantees). The adversary observes all
transcripts. Upon corruption, it learns the corrupted party's current state
(set $X_\tau$, key $K_\tau$, IBLT, VOLE/OKVS state). Forward privacy requires
learning \emph{nothing} about sets from epochs $t < \tau$ beyond what union
outputs already reveal. Post-compromise security is not required in the main
protocol; a self-healing extension appears in Appendix~\ref{sec:selfheal}.

The formal ideal functionality $\cF_{\mathsf{updPSU}}$
(Fig.~\ref{fig:f-fpsu}) and UC definitions appear in
Sections~\ref{sec:protocol} and~\ref{sec:security}.
$\cF_{\mathsf{updPSU}}$ receives deltas, computes the updated union and
intersection, and returns the output corresponding to a mode parameter
($\mathsf{PSU}, \mathsf{PSI}, \mathsf{PSU}\text{-}\mathsf{CA},
\mathsf{PSI}\text{-}\mathsf{CA}$). Theorem~\ref{thm:lower} proves the
leaked delta counts are unavoidable.

% ============================================================
\subsection{Leakage Lower Bound}\label{sec:lower}

\begin{theorem}[Efficiency-Leakage Trade-off]\label{thm:lower}
Any correct, semi-honest secure, 1-updatable PSU protocol necessarily reveals
$(|\Delta^+_0|, |\Delta^-_0|, |\Delta^+_1|, |\Delta^-_1|)$. Any protocol
hiding these counts must have $|\mathsf{trans}| = \Omega(N\sigma)$---defeating
updatability.
\end{theorem}

F-PSU leaks these four counts through OT query counts and delta IBLT sizes.
The proof uses a Shannon source-coding argument. F-PSU leaks exactly these
four counts; by Theorem~\ref{thm:lower}, it is \textbf{strictly
leakage-optimal}. The full proof appears in Appendix~\ref{sec:lower-proof}.

\subsection{Protocol Overview}

We sketch one incremental epoch (the full specification is
Figure~\ref{fig:protocol}). After Epoch~0 establishes the shared-element
IBLT, each subsequent epoch proceeds in three stages.

\smallskip\noindent\textbf{Additions (MiniPeel).} Each party encodes its
new elements into a compact IBLT. A fresh UnionPeel on these IBLTs
recovers the union of additions deterministically, with zero cascade.

\smallskip\noindent\textbf{Deletions (OT-masked).} Before modifying the
IBLT, each party queries the other via batched OPRF plus one deletion-intent
bit. Only when both parties delete the same shared element (rare) is the
IBLT modified; otherwise deletions are handled purely locally or via $k$
cell updates, without cascade.

\smallskip\noindent\textbf{Output and forward privacy.} The new union is
assembled from the old union and classified deltas, then the PRF key is
evolved, the IBLT encrypted, and the old key erased.

The per-epoch cost is $O(|\Delta|)$. The detailed protocol follows in
Sections~\ref{sec:ideal}--\ref{sec:complexity}.

\subsection{Ideal Functionality}\label{sec:ideal}

	In the UC framework~\cite{Can01}, an ideal functionality $\cF_{f}$ is
	parameterized by the function $f$ it computes. We define
	$\cF_{\mathsf{updPSU}}$ for forward-private updatable Private Set Union
	(Fig.~\ref{fig:f-fpsu}).

	\begin{center}
	\fbox{\begin{minipage}{0.95\columnwidth}
	\smallskip
	\noindent\textbf{Functionality $\cF_{\mathsf{updPSU}}$}

	\noindent\textbf{Parameters:} Two parties $P_0, P_1$, security parameter
	$\kappa$, element domain $\U$, and max set size $N$. The functionality
	maintains set state $(X_0, X_1)$ and epoch counter $t$ (initially $0$).

	\smallskip\noindent\textbf{Behavior:}
	\begin{enumerate}
	\item Upon receiving $(\textsf{Update}, \Delta^+_b, \Delta^-_b)$ from $P_b$:
	update $X_b \leftarrow (X_b \cup \Delta^+_b) \setminus \Delta^-_b$.
	\item Once both parties have submitted their updates, set $t \leftarrow t+1$,
	compute $\mathsf{out} \leftarrow X_0 \cup X_1$, and send $\mathsf{out}$
	to both parties.
	\item Upon adaptive corruption of $P_b$ at epoch $\tau$: send
	$X_b^{(\tau)}$ to $\cS$. Historical states $\{X_b^{(t)}\}_{t < \tau}$.
	\end{enumerate}
	\smallskip\noindent\textbf{Leakage:} At step~2, $\cL = (|\Delta^+_0|,
	|\Delta^-_0|, |\Delta^+_1|, |\Delta^-_1|)$ to the simulator.
	\smallskip
	\end{minipage}}
\captionof{figure}{The ideal functionality $\cF_{\mathsf{updPSU}}$ for forward-private updatable Private Set Union.}
\label{fig:f-fpsu}

\smallskip\noindent\textbf{Remark (PSI and cardinality as a byproduct).}
The UnionPeel sub-protocol naturally classifies each peeled element as unique
to $P_0$, unique to $P_1$, or shared. The same protocol execution therefore
also yields the intersection (shared elements) and the cardinalities of both
the union and the intersection, at zero additional communication or computation.
The security proof is unchanged since the protocol transcript is independent
of which output elements the parties retain. We focus on $\cF_{\mathsf{updPSU}}$
throughout; the PSI and CA variants follow identically.
\end{center}

	We state the main security theorem for $\cF_{\mathsf{updPSU}}$;
	the proofs for the PSI and cardinality variants follow identically
	(Remark above), since the protocol transcript is independent of which
	peeled elements the parties retain.

\subsection{Extension: PSI and Cardinality from Pure Cells}
\label{sec:psi-ext}

The $\uPeel$ function (Definition~\ref{def:upeel}) naturally reveals whether a
peeled element belongs to the intersection. When $\uPeel$ peels a cell
$(i,j)$, it distinguishes three cases: (1)~$\mathsf{cnt}_0[i][j]=1,
\mathsf{cnt}_1[i][j]=0$---the element is unique to $P_0$;
(2)~$\mathsf{cnt}_0[i][j]=0, \mathsf{cnt}_1[i][j]=1$---unique to $P_1$;
(3)~$\mathsf{cnt}_0[i][j]=1, \mathsf{cnt}_1[i][j]=1, \mathsf{sum}_0=
\mathsf{sum}_1$---the element belongs to \textbf{both} parties, i.e., the
intersection $X_0 \cap X_1$. Case~(3) occurs because each party independently
encodes its elements into its own IBLT; when the same element $x$ is encoded
by both, it lands in the same cell $(i,j)$ with identical sum $x$, and
$\uPeel$ detects the match.

Consequently, a single execution of the UnionPeel sub-protocol simultaneously
recovers the union (all peeled elements) and the intersection (Case~3
elements). This yields four output types at zero additional communication:
\begin{itemize}
\item \textbf{PSU:} output all peeled elements ($X_0 \cup X_1$).
\item \textbf{PSI:} output only Case~3 elements ($X_0 \cap X_1$).
\item \textbf{PSU-CA:} output the count $|X_0 \cup X_1|$.
\item \textbf{PSI-CA:} output the count $|X_0 \cap X_1|$.
\end{itemize}

The mode is selected via a single bit in the protocol initialization and
requires no change to the protocol messages, the OT/OPRF interaction, the
IBLT encoding, or the peeling logic. The only modification is that the
peeling output carries a 2-bit tag per element identifying the case, and the
parties filter output accordingly.

Security inherits directly from Theorem~\ref{thm:main}: the simulator for
$\Pi_{\mathsf{fpsu}}$ already extracts the corrupted party's effective input
and submits it to the ideal functionality. The same simulator, paired with
$\cF_{\mathsf{updPSU}}$ (which differs only in the output sent to the parties),
produces an indistinguishable view. The per-mode leakage is identical:
$X_\cup$ reveals $Y \setminus X$ in all modes; $X_\cap$ additionally reveals
element equality between the two sets; cardinalities reveal only the sizes.
The adaptive security argument of Section~\ref{sec:security} applies
unchanged since the protocol transcript is independent of the output mode.

\subsection{Initialization}
\label{sec:init}

At setup, $P_0$ and $P_1$ (denoted Client $C$ and Server $S$ in pseudocode
perform one-time VOLE setup (RR22~\cite{RR22}) reused across all epochs.
Each party initializes an empty IBLT of capacity $(m, k)$ and samples an
independent 256-bit PRF key.

\subsection{Epoch Protocol}\label{sec:epoch-protocol}

	We denote the two parties as Client $C$ and Server $S$, each initialized with
	the one-time setup of Section~\ref{sec:init}.  Let
	$\mathsf{prm}_t = (M, k, \ell_t, H)$ denote IBLT parameters chosen such that
	$\List$ fails with probability $\leq 2^{-\lambda}$ for the current set sizes
	(Theorem~\ref{thm:iblt}), with expansion $e = k\ell_t / (|X_t| + |Y_t|) \in
	[1.5, 4.5]$.  The per-epoch protocol is specified in
	Figure~\ref{fig:protocol}; we elaborate on the key design decisions below.

	\smallskip\noindent\textbf{OT-masked deletion.}
	\smallskip\noindent\textbf{Mini UnionPeel for additions.}
	Additions are encoded into a separate compact IBLT ($\mathsf{IBLT}^+_C$,
	$\mathsf{IBLT}^+_S$) with $m_{\Delta} = c_k \cdot k \cdot \max(|\Delta^+_C|,
	|\Delta^+_S|)$ cells, \emph{without} merging into the main IBLT.  A fresh
	$\Pi_{\mathsf{UnionPeel}}$ on this compact IBLT recovers the union of
	additions deterministically (zero cascade, as the IBLT has no prior peeling
	state).  Recovered elements $V_{\Delta^+}$ are then merged locally.  This
	separation avoids cross-contamination: without it, deleting $a$ and adding
	$b$ that share a hash cell would produce a non-singleton with sum $b-a$.

	\smallskip\noindent\textbf{OT-masked deletion.}
	The critical insight of F-PSU is that deletions need not trigger IBLT
	cascade. Before modifying any IBLT cell, each party runs a 2-bit OT-based
	membership check for every deletion candidate (Step~5, Fig.~\ref{fig:protocol}).
	The other party replies with (i)~whether it holds the element, and
	(ii)~whether it also intends to delete it. Three cases arise:
	\begin{itemize}
	  \item \textbf{Case~1 (unique):} The other party does not hold $x$. Pure
	    local deletion---no IBLT modification, no communication beyond the OT check.
	  \item \textbf{Case~2 (one-sided shared):} The other party holds $x$ and
	    retains it. $x$ remains in the union. The deleting party subtracts $x$
	    from its own IBLT ($k$ cell updates, no cascade).
	  \item \textbf{Case~3 (two-sided shared):} Both parties independently
	    delete $x$. This is the only case requiring IBLT subtraction from both
	    parties.
	\end{itemize}
	By avoiding cascade for Cases~1 and~2, the chain reaction that would
	otherwise peel through all remaining shared elements is eliminated.

	\smallskip\noindent\textbf{Output construction.}
	The new union is $X_{\cup, t} = (X_{t-1} \setminus \Delta^-_{\mathsf{Case1}})
	\cup \Delta^+_C \cup \Delta^+_S \cup V_{\mathsf{main}}$, where
	$\Delta^-_{\mathsf{Case1}}$ are unique deletions and $V_{\mathsf{main}}$
	is the (usually empty) set of elements resolved from Case~3 IBLT cells
	(Step~6).  If the output type is $\mathsf{PSI}$, only Case~3 elements of
	the original $\uPeel$ definition (equal sums in both IBLTs) are retained.

	\smallskip\noindent\textbf{Key evolution and forward privacy.}
	Round~3 performs a forward-secure key update: each party derives
	$K^b_t \leftarrow \PRF(K^b_{t-1}, \text{``evolve-}b\text{''})$ (party-specific
	domain separator), encrypts its IBLT under $K^b_t$, writes the
	ciphertext, and \emph{securely erases} $K^b_{t-1}$.  By
	Lemma~\ref{lem:prf-ow}, an adversary compromising at epoch $\tau$ cannot
	recover $K^b_{t<\tau}$ or decrypt prior-epoch IBLT ciphertexts.

	\begin{figure*}[p]
\fbox{\footnotesize\begin{minipage}{\dimexpr\textwidth-2\fboxsep-2\fboxrule\relax}
\smallskip
\noindent\textbf{Protocol $\Pi_{\mathsf{fpsu}}$ (Epoch $t$)}

\smallskip\noindent\textbf{Input of $C$:}
Encrypted IBLT $\mathsf{CT}_C$, set $X_{t-1}$, key $K^C_{t-1}$,
deltas $\Delta^+_C, \Delta^-_C$.

\noindent\textbf{Input of $S$:}
Encrypted IBLT $\mathsf{CT}_S$, set $Y_{t-1}$, key $K^S_{t-1}$,
deltas $\Delta^+_S, \Delta^-_S$.

\noindent\textbf{Common:}
IBLT parameters $\mathsf{prm}_t = (M, k, \ell_t, H)$ with $k=5$, $e \in [1.5,4.5]$;
VOLE state $\sigma_{\mathsf{VOLE}}$ (setup once, reused across all epochs).

\noindent\textbf{Output:}
Both parties output $X_{\cup,t}$. (By filtering in Step~8, also yields
PSI, PSU-CA, PSI-CA.)

\smallskip\noindent\textbf{Protocol:}
\begin{enumerate}
\item $C$ decrypts $\mathsf{IBLT}_C \leftarrow \SE.\mathsf{Dec}(K^C_{t-1}, \mathsf{CT}_C)$;
      $S$ decrypts $\mathsf{IBLT}_S \leftarrow \SE.\mathsf{Dec}(K^S_{t-1}, \mathsf{CT}_S)$.

\item $C$ encodes additions: $\mathsf{IBLT}^+_C \leftarrow \Encode(\Delta^+_C)$;
      $S$ encodes additions: $\mathsf{IBLT}^+_S \leftarrow \Encode(\Delta^+_S)$,
      each with $m_{\Delta} = c_k \cdot k \cdot \max(|\Delta^+_C|, |\Delta^+_S|)$ cells.
      \emph{(Do not merge into main IBLT yet.)}
\end{enumerate}
\begin{enumerate}\setcounter{enumi}{2}
\item \textbf{[Mini UnionPeel --- additions.]}
      $C$ and $S$ run $\Pi_{\mathsf{UnionPeel}}(\mathsf{IBLT}^+_C, \mathsf{IBLT}^+_S)$
      (Fig.~\ref{fig:unionpeel}), scanning all $m_{\Delta}$ cells. The addition IBLT
      is fresh --- zero cascade, deterministic output $V_{\Delta^+}$.

\item \textbf{[Merge additions.]} [local]
      Shared additions go to both IBLTs; unique additions go only to the owning party's IBLT.
\end{enumerate}
\begin{enumerate}\setcounter{enumi}{4}
\item \textbf{[OT-masked deletion.]}
      For each $x \in \Delta^-_C$, $C$ and $S$ run a batched OPRF query. $S$ replies with $F(k,x)$ and two bits:
      $S$ replies with $(b_1, b_2)$ where $b_1 = \mathbf{1}[x \in Y_{t-1}]$ and
      $b_2 = \mathbf{1}[x \in \Delta^-_S]$. Three cases:
      \begin{itemize}
        \item[Case~1:] $b_1=0$: $x$ unique to $C$, local deletion only.
        \item[Case~2:] $b_1=1, b_2=0$: $x$ shared, $S$ keeps it ---
          $\mathsf{IBLT}_C \leftarrow \mathsf{IBLT}_C \ominus \Encode(\{x\})$ (local, $k$ cells, no cascade).
        \item[Case~3:] $b_1=1, b_2=1$: both delete $x$ --- $\mathsf{IBLT}_C \leftarrow \mathsf{IBLT}_C \ominus \Encode(\{x\})$, record cell to $Q_{\mathsf{del}}$.
      \end{itemize}
      $S$ performs the symmetric check for each $y \in \Delta^-_S$.
      Most deletions are Case~1 or Case~2
      eliminating cascade propagation.

\item \textbf{[Resolve Case~3 deletions.]}
      If $Q_{\mathsf{del}} \neq \emptyset$: run $\Pi_{\mathsf{UnionPeel}}$ on $Q_{\mathsf{del}}$.
      By Lemma~\ref{lem:cascade}, $\leq 2k|\Delta^-_{\mathsf{Case3}}|$ cells examined
      w.h.p.

\item \textbf{[Transmit recovered elements.]} [$S \to C$]
      $S$ sends $V_{\mathsf{main}}$ (resolved elements from Steps~4--6, if any) to $C$.
\end{enumerate}

\begin{enumerate}\setcounter{enumi}{7}
\item \textbf{[Output.]}
      $X_{\cup, t} \leftarrow (X_{t-1} \setminus \Delta^-_{\mathsf{Case1}}) \cup
      \Delta^+_C \cup \Delta^+_S \cup V_{\mathsf{main}}$.
      Both parties output $X_{\cup, t}$.

\item \textbf{[Key evolution.]}
      $C$: $K^C_t \leftarrow \PRF(K^C_{t-1}, \text{``evolve-C''})$;\;
      $S$: $K^S_t \leftarrow \PRF(K^S_{t-1}, \text{``evolve-S''})$.

\item \textbf{[Re-encrypt and erase.]}
      $C$: $\mathsf{CT}_C \leftarrow \SE.\mathsf{Enc}(K^C_t, \mathsf{IBLT}_C)$,
      $\mathsf{secure\_erase}(K^C_{t-1})$;
      $S$: $\mathsf{CT}_S \leftarrow \SE.\mathsf{Enc}(K^S_t, \mathsf{IBLT}_S)$,
      $\mathsf{secure\_erase}(K^S_{t-1})$.
\end{enumerate}
\smallskip
\end{minipage}}
\caption{The F-PSU epoch protocol. Steps~2--4: Mini UnionPeel for additions
(deterministic, zero cascade). Steps~5--7: OT-masked deletion with three-case
classification.
Steps~8--10: output, key evolution, and secure erasure. Per-epoch cells
scanned: $O(|\Delta^+|)$ (MiniPeel) $+$ $O(|\Delta^-_{\mathsf{Case3}}|)$
(Case~3 cleanup). Forward privacy overhead: $<1$~ms.}
\label{fig:protocol}
\end{figure*}

\subsection{Complexity}\label{sec:complexity}

\textbf{Unbalanced settings.} F-PSU operates identically in balanced and
unbalanced scenarios because IBLT encoding complexity depends only on the
\emph{local} set size, not the peer's set size. When $P_A$ holds $2^{10}$
elements and $P_B$ holds $2^{20}$, each party encodes its own IBLT with cost
proportional to $|\text{local set}|$, independent of the other. The delta
IBLT mechanism is similarly local: each party applies $\Delta$ only to its
own IBLT. The UnionPeel sub-protocol treats both IBLTs symmetrically,
peeling from the virtual union regardless of size disparity. No protocol
modification is required. The asymptotic complexity is $O(n\sigma)$ for the
union output and $O(d(\sigma + \kappa))$ for the per-epoch update, where
$n = |X| + |Y|$ and $d = |\Delta|$.

Let $d = |\Delta| = d^+ + d^-$, $n = |X| + |Y|$.
\begin{itemize}
\item \textbf{OT-masked deletion:} $O(d^-)$ OT membership checks (2-bit OT per
element).

\item \textbf{Mini UnionPeel:} $O(c_k \cdot k \cdot d^+ \cdot (\sigma + \kappa))$ bits,
deterministic, zero cascade.
\item \textbf{Output:} $n\sigma$ bits (inherent to any PSU).
\item \textbf{Forward privacy:} one PRF + one symmetric encryption per epoch ($<1$~ms).
\end{itemize}
Total: $O(d(\sigma + \kappa) + n\sigma)$, dominated by MiniPeel for additions
and OT checks for deletions. For $d/n \approx 0.1\%$, the non-output
component is $\approx 1000\times$ smaller than static PSU~\cite{PT26}.

\begin{table}[ht]
\centering
\footnotesize
\caption{Asymptotic comparison of PSU protocols ($d = |\Delta|$, $n = |X|+|Y|$).}
\label{tab:complexity}
\resizebox{\columnwidth}{!}{%
\begin{tabular}{lccc}
\toprule
\textbf{Protocol} & \textbf{Communication} & \textbf{Computation} & \textbf{Bal./Unbal.} \\
\midrule
KRT19~\cite{KRT19} & $O(n\sigma\kappa)$ & $O(n\log n)$ OT & Balanced \\
PT26~\cite{PT26} & $O(n(\sigma+\kappa))$ & $O(n)$ OT+hash & Balanced \\
ePSU~\cite{Tu+25} & $O(n\sigma)$ & $O(n)$ symm. & Both \\
\textbf{F-PSU} & $\mathbf{O(d(\sigma+\kappa)+n\sigma)}$ & $\mathbf{O(d)}$ OT$+\mathbf{O(n)}$ ChCha & \textbf{Both} \\
\bottomrule
\end{tabular}%
}
\end{table}

\begin{lemma}[Cascade Bound]\label{lem:cascade}
With IBLT parameters ensuring 2-core emptiness ($\Pr \geq 1 - 2^{-\lambda}$),
peeling from $k d^-$ deletion-modified cells examines at most $(1 + C_0) \cdot
k d^-$ total cells, where $C_0 = O(k^2 d^- / m)$. For $d^- \ll m/k$, the
cascade is $\leq 2k d^-$ with overwhelming probability.
\end{lemma}

\subsection{Communication Reduction}

Table~\ref{tab:main-perf} compares per-epoch IBLT data exchange against static
re-encoding. Communication reduction scales as $\Theta(n/|\Delta|)$,
ranging from $29\times$ ($n{=}2^{14}, |\Delta|{=}1024$) to
$23{,}756\times$ ($n{=}2^{20}, |\Delta|{=}64$). This reflects the core
property of OT-masked deletion: per-epoch cell scanning is dominated by
MiniPeel ($O(|\Delta^+|)$ cells), with deletion cascade eliminated by
offloading membership checks to OT. OT/OPRF communication scales
proportionally with cell count, so the same reduction factor applies to the
secure computation phase.

\subsection{Incremental Peeling and Parameter Growth}

For $n = 2^{16}$, $|\Delta| = 2^{12}$: incremental peeling examines 32,768
cells vs.\ 197,000 for full scan ($6\times$ reduction). Per-epoch protocol:
$20.8$~ms ($n = 2^{14}$, $|\Delta| = 2^{12}$), $97.8$~ms ($n = 2^{16}$,
$|\Delta| = 2^{14}$). Forward privacy overhead: $<1$~ms.

\textbf{Amortized parameter growth.} 200 epochs, $n_0 = 10^4$, 100 elements
per epoch: only \textbf{1} full re-encode triggered (epoch~51).

% ============================================================